\section{Problem Formulation}
\label{sec:problem}

\subsection{Setting}

A multi-agent LLM coding session is a sequence of \emph{turns}
\(\turn_1, \turn_2, \ldots, \turn_N\) on the orchestrator's main
thread, interleaved with subagent dispatches.  Each turn is processed
on a model tier \(\tier \in \tierset = \{\text{haiku}, \text{sonnet},
\text{opus}\}\), ordered by capability and price.  Let
\(\alpha_{\tier}\) and \(\beta_{\tier}\) denote the input- and
output-token rates (USD per token) for tier \(\tier\).  For a turn
\(\turn_i\) with context length \(|\text{ctx}(\turn_i)|\) input tokens
and \(|\text{out}(\turn_i)|\) output tokens, the raw per-turn cost is:

\begin{equation}
  \cost(\turn_i, \tier) =
    \alpha_{\tier} \cdot |\text{ctx}(\turn_i)| \;+\;
    \beta_{\tier} \cdot |\text{out}(\turn_i)|.
  \label{eq:perturn}
\end{equation}

Prefix caching introduces a discount \(\rho \in [0,1]\) on the cached
fraction of the input, so the effective input cost becomes
\(\alpha_{\tier}\bigl((1-\rho)\,|\text{ctx}(\turn_i)| +
\rho\,|\text{ctx}(\turn_i)| \cdot r_{\text{cache}}\bigr)\), where
\(r_{\text{cache}} = 0.10\) is Anthropic's cache-read rate
relative to the base rate~\cite{anthropic-cache}.  For notational
convenience we fold this discount into \(\alpha_{\tier}\) and write
the effective rate \(\hat\alpha_{\tier}(\rho)\).

\subsection{Empirical motivation}

We measure \blackrim's cost split as 99.4\,\% main-thread (Gestalt,
opus), 0.6\,\% subagent dispatch over a 7-day production window
(Table~\ref{tab:cost-split}, §\ref{sec:eval}).  The dispatched
subagents are already routed by an advisor that applies a 16-rule
role-\(\times\)-signal cascade~\cite{moa1landscape}; the 0.6\,\%
residual reflects a well-tuned dispatch-side system.  The orchestrator
itself is the dominant cost lever.

\subsection{Per-turn routing decision}

For each main-thread turn \(\turn_i\), choose tier \(\tier_i\) to
minimise expected cost subject to a quality-preservation constraint:

\begin{equation}
  \tier_i^{*} = \arg\min_{\tier \in \tierset}\; \cost(\turn_i, \tier)
  \quad \text{s.t.} \quad
  \mathbb{E}[\quality(\turn_i, \tier)] \;\geq\; \quality_0 - \epsilon,
  \label{eq:routing}
\end{equation}

where \(\quality_0 = \mathbb{E}[\quality(\turn_i, \text{opus})]\) is
the opus-tier quality baseline and \(\epsilon\) is the
per-turn quality-loss tolerance.  In practice \(\quality\) is
observable post-hoc through three channels~\cite{moa1landscape}:
structured-output correctness, tool-use success, and eval-suite scores
when suites exist.  The tier-down condition
\(\cost(\turn_i, \tier_{\text{low}}) \leq \cost(\turn_i, \tier_{\text{high}})\)
is always satisfied for any \(\tier_{\text{low}} \prec \tier_{\text{high}}\)
in the price ordering; the binding constraint is quality.

\subsection{Plan-cache decision}

For each subagent dispatch, the orchestrator emits a \emph{plan}
with semantic signature \(\plan\) (§\ref{sec:plancache}).  The
plan-cache decision is whether to reuse a previously computed plan
output or incur the full plan-composition cost:

\begin{equation}
  \cost_{\text{cached}}(\plan) =
    \hitrate(\plan) \cdot \cost_{\text{lookup}} \;+\;
    \bigl(1 - \hitrate(\plan)\bigr) \cdot \cost_{\text{compose}}(\plan),
  \label{eq:plancache}
\end{equation}

where \(\hitrate(\plan)\) is the hit rate for plans whose signature
matches \(\plan\) at threshold, \(\cost_{\text{lookup}}\) is the
cache-fetch cost (an O(1) hash lookup plus a short
embedding pass; negligible token cost), and
\(\cost_{\text{compose}}(\plan)\) is the standard plan-composition
cost (retrieval, similar-cases lookup, briefing assembly via
\texttt{gt context budget}).  The cache reuse condition is:
\(\operatorname{sim}(\plan_{\text{query}}, \plan_{\text{cached}}) > \theta_{\text{hit}}\)
and \(\hitrate(\plan) > \theta_{\text{conf}}\), where
\(\theta_{\text{hit}}\) and \(\theta_{\text{conf}}\) are threshold
parameters calibrated on the evaluation set.

\subsection{Total session cost}

Over \(N\) main-thread turns and \(M\) dispatches, the session cost is:

\begin{equation}
  \cost_{\text{session}} =
    \underbrace{\sum_{i=1}^{N} \cost(\turn_i, \tier_i)}_{\text{main-thread routing}} \;+\;
    \underbrace{\sum_{j=1}^{M} \cost_{\text{cached}}(\plan_j)}_{\text{plan-cache dispatch}}.
  \label{eq:session}
\end{equation}

The two terms are independently optimised: routing chooses
\(\tier_i\) per turn on the main thread, while plan caching reduces
the expected-compute term within each dispatch.

\subsection{Composition and independence}

Under the assumption that routing decisions are independent of plan
signatures (i.e., tier choice does not systematically correlate with
plan structure), the two mechanisms contribute additively to cost
reduction.  In practice the composition is slightly superadditive:
plan-cache hits bypass the briefing-composition LLM call entirely,
so they avoid both the creation-side cache-write cost and any
cache-miss penalty.  This is distinct from prefix-cache savings, which
operate on the static prefix of already-launched LLM calls.  We
discuss the composition in §\ref{sec:plancache} and measure it
empirically in §\ref{sec:eval}.
